\section{The energy method with prescribed columns}\label{sec:energy}

This section recalls OpenAI's energy method, extends it to determinants with prescribed columns, and
proves Theorem~\ref{thm:real} assuming the bound of Section~\ref{sec:signed}. Throughout, $\alpha_1=\AlphaOne$.

\subsection{OpenAI's energy}\label{ssec:energy}
By Andr\'eief's identity, $\det\cA_N$ is an integral over $n$ points $x_i\in(-1,1)$ and $n$ points
$s_j\in(0,1)$. OpenAI bounds the integrand by $e^{n^2E}$, where $E$ is a discrete logarithmic energy of the
points, written in the coordinate $x$ with $t=2x/(1+x^2)$ \cite[\S5]{OpenAI2026}. It has the form
\begin{multline*}
 n^2E=\sum_iV_1(x_i)+\sum_jV_2(s_j)+\sum_{i<i'}\bigl(\log|x_i-x_{i'}|+\log(1-x_ix_{i'})\bigr)\\
 +\sum_{j<j'}\log|s_j-s_{j'}|-\sum_{i,j}\log(1-2x_is_j+x_i^2),
\end{multline*}
schematically, with external fields $V_1,V_2$ of order $n$. The first two sums are the
external fields. The next two are the logarithmic interactions among points of the same kind, and the last
is the interaction between the two kinds. There are two cases, according to which of OpenAI's two bounds for
the determinant of the row functions is used \cite[Props.~5.2, 5.3]{OpenAI2026}. The case with doubled interaction is the one in which OpenAI's condition
$\sum_i(1-x_i^2)/(1+x_i^2)\le n-1-2g$ \cite[Eq.~(67)]{OpenAI2026} holds, and the constrained case is its complement. In the \emph{constrained
case} the row field acquires a term with a multiplier $\lambda_1$ for the constraint that defines the case.
In the \emph{case with doubled interaction} the interaction among the $x_i$ has weight $2$ instead of $1$. This case also bounds a determinant
$\det\mathsf L$ of size $n$ by $e^{O(n\log n)}$, which is the factor that Section~\ref{sec:signed} improves.

The energy is bounded from above by duality. Given trial sequences $u=(u_k)$ and $v=(v_k)$, one writes the
interaction terms as quadratic forms in the moments $\frac1n\sum_iT_k(x_i)$ and $\frac1n\sum_js_j^k$. Then one
completes the squares against $u$ and $v$. What remains is the sum of the norms $\|u\|^2=\sum_ku_k^2/k$ and
$\|v\|^2$ and, for each point, a function of one variable. With $\mathsf T(u,x)=\sum_ku_kT_k(x)/k$ and
$\mathsf S(v,x)=\sum_kv_kx^k/k$, these functions are, in the case with doubled interaction,
\begin{align*}
 X(x)&=\alpha_A\log|x|+\beta_h\log(1-x)-\gamma_H\log(1+x^2)-4\mathsf T(u,x)-\mathsf S(v,x),\\
 Y(s)&=\tfrac{b_0}{n_0}\log s+\tfrac{q_0}{n_0}\log(1-s)+2\mathsf T(v,s).
\end{align*}
In the constrained case, with a trial sequence $u'$ in place of $u$, they are
\begin{multline*}
 X_1(x)=\alpha_A\log|x|+\beta_h\log(1-x)-(\gamma_H-\tfrac12)\log(1+x^2)\\
 +\lambda_1\bigl(\tfrac{2g_0}{n_0}-\tfrac{2x^2}{1+x^2}\bigr)-2\mathsf T(u',x)-\mathsf S(v,x),
\end{multline*}
with the same $Y$. If $X\le M_X$ and $Y\le M_Y$ on the whole domains, the energy of every configuration is at
most a constant determined by $M_X$, $M_Y$ and the norms. In the language of potential theory the extremal
configuration is described by an equilibrium problem for two interacting measures in an external field
\cite{SaffTotik}. The best trial sequences are the moments of its solution, as Carreira observed
\cite{Carreira2026}. Any trial sequences give a valid upper bound.

\subsection{Prescribed points}\label{ssec:prescribed}
When $m$ of the $n$ columns are point values, only $n-m$ points of each kind are integrated over, and the
energy is still normalized by $n$. The prescribed points enter through the row determinant, now evaluated
at the free points together with the prescribed ones. There are two kinds of prescribed points: the
quadrature nodes, which lie in the set $E_x$ of Section~\ref{sec:real}, and points $y\in(0,1)$. Let $\alpha=m/n$ be the fraction of
prescribed columns and $\alpha_2$ the fraction of prescribed points of the second kind. Each missing free point
changes the energy bound by a bounded amount. The total change is linear in $m$, and its slope $c_2$ enters
$\rho$.

Every step of OpenAI's argument remains valid uniformly in the number of free points. The regularization of
the logarithmic interactions, the diagonal corrections and the completion of squares only involve sums over
the points present. For the prescribed points of the first kind the relevant field is
\[
 F_1(x)=-\gamma_d\log(1+x^2)-2\mathsf T(u,x),\qquad x\in E_x .
\]
For those of the second kind it is
\[
 F_2(y)=-\gamma_d\log(1+y^2)+\frac{n_0-2g_0-1}{n_0}\log y-u_1y,\qquad0<y<1,
\]
where $u_1$ is the first coefficient of $u$. The coefficient $(n_0-2g_0-1)/n_0$ is a lower bound, valid for all
$N\ge1$, of $(n-1-2g)/n$. By Lemma~\ref{lem:binomvand} the generalized Vandermonde determinant contributes only
the factor $\prod_{i<l}|s_i-s_l|$ to the energy, and the binomial factors are assigned to the columns. The
energy bound then reduces to a quadratic inequality in $\alpha$ and $\alpha_2$.

\begin{lemma}[energy inequality with prescribed points]\label{lem:energy-ineq}
Let $\ell=\log2$, $\alpha_1>0$ and $u\ge0$, and let $\mathcal N,\mathcal F,\mathcal R,\mathcal F_2$ and $c_2$ be real
numbers with
\begin{gather*}
 \mathcal N+\mathcal F+U_0-a_0\ell\le u,\qquad \ell-\mathcal R+\mathcal F_2\le0,\\
 (a_1-1)\ell-\mathcal F+\mathcal R-c_2-\ell\alpha_1+\frac u{\alpha_1}\le0 .
\end{gather*}
Then for all $\alpha\ge\alpha_1$ and $\alpha_2\ge0$
\begin{multline*}
 \bigl(-(a_0-\tfrac12)+(a_1-1)\alpha-\alpha^2+\alpha_2-\alpha_2^2\bigr)\ell+\mathcal N+(1-\alpha)\mathcal F\\
 +(\alpha-\alpha_2)\mathcal R+\alpha_2\mathcal F_2\le-U_0+\tfrac12\ell+c_2\alpha .
\end{multline*}
\end{lemma}

\begin{proof}
Put $s=(a_1-1)\ell-\mathcal F+\mathcal R-c_2$. The left side minus the right side equals
\[
 (\mathcal N+\mathcal F+U_0-a_0\ell)+\alpha s-\ell\alpha^2+\alpha_2(\ell-\mathcal R+\mathcal F_2)-\ell\alpha_2^2 .
\]
The last two terms are nonpositive by the second hypothesis. The first three are at most
$\alpha(u/\alpha+s-\ell\alpha)$. The function $u/\alpha-\ell\alpha$ decreases in $\alpha$, so this is at most
$\alpha(u/\alpha_1+s-\ell\alpha_1)\le0$.
\end{proof}

The lemma is applied with $\mathcal N$ a bound for $2\|u\|^2+\frac12\|v\|^2$, $\mathcal F=M_X+M_Y$,
$\mathcal R$ a bound for the supremum of $F_1$ and $\mathcal F_2$ a bound for the supremum of $F_2$. The
analogue for the constrained case has no prescribed points of the second kind. Its polynomial in $\alpha$ is
\[
 \bigl(-(a_0-\tfrac12)+(a_1-1)\alpha-\alpha^2\bigr)\ell+n_1+(1-\alpha)\mathcal F_1+\alpha\lambda_1\tfrac{2g_0}{n_0}
 \le-U_0+\tfrac12\ell+c_2\alpha,
\]
with $n_1$ a bound for the norms of the trial sequences and $\mathcal F_1=M_{X_1}+M_Y$. The coefficient of $\alpha^2$ is negative, so this holds for all real $\alpha$ as soon as the discriminant of the
difference of the two sides is nonpositive.

For our parameters the trial sequences of Appendix~\ref{app:field} give the hypotheses of
Lemma~\ref{lem:energy-ineq} with $\alpha_1=\AlphaOne$, $u=0.0068682$ and $c_2=\Ctwo$. They also give the
inequality of the constrained case with the same $U_0$ and $c_2$.

\emph{The squares that remain.} Section~\ref{sec:signed} uses what the completion of squares leaves, so we record it for the case with doubled
interaction. Let $\mu_F$, $\mu_M$ and $\mu_S$ be the normalized counting measures of the free points, the prescribed points of the first kind and the
prescribed points of the second kind. Their total masses are $1-\alpha$, $\alpha-\alpha_2$ and $\alpha_2$. Put $W=\mu_F+\mu_M$,
$A_k(\mu)=\int T_k\,d\mu$ and $\langle a,b\rangle=\sum_ka_kb_k/k$, so that $\|a\|^2=\langle a,a\rangle$. Three products carry the interactions
in the variable $x$. The Vandermonde product of the free points comes from the Cauchy determinant. The row determinant gives the Vandermonde product of the
free points and the prescribed points of the first kind, and the Vandermonde product of the points $y$ of the second kind. These last points
enter the row determinant through $1/y$, so it also gives the cross factors $|w-1/y|=(1-wy)/y$, for $w$ a point of $W$.
Since $-\log|x-x'|=\log2+2\sum_kT_k(x)T_k(x')/k$ and $\log(1-wy)\le-wy$, the part of $E$ quadratic in the Chebyshev moments is, apart from the
diagonal terms and the multiples of $\log2$, at most
\[
 -\|A(\mu_F)\|^2-\|A(W)\|^2-\|A(\mu_S)\|^2-A_1(W)A_1(\mu_S).
\]
For any measure $\mu$ we have $-\|A(\mu)\|^2=\|u\|^2-2\langle u,A(\mu)\rangle-\|A(\mu)-u\|^2$. Using this for $\mu_F$ and $W$, and writing
$A_1(W)=u_1+(A_1(W)-u_1)$, we find that the display equals
\[
 2\|u\|^2-2\langle u,A(2\mu_F+\mu_M)\rangle-u_1A_1(\mu_S)-\mathcal Q,
\]
where
\[
 \mathcal Q=\|A(\mu_F)-u\|^2+\|A(W)-u\|^2+\|A(\mu_S)\|^2+(A_1(W)-u_1)A_1(\mu_S).
\]
The first three terms are the term $2\|u\|^2$ of $\mathcal N$ and the terms of $X$, $F_1$ and $F_2$ that contain $u$. So the energy bound obtained from
Lemma~\ref{lem:energy-ineq} improves by $\mathcal Q$. To bound $\mathcal Q$ from below, let $e_1$ be the first
unit vector and put $a=A(\mu_F)-u$ and $b=A(W)-u+\frac12A_1(\mu_S)e_1$. Since $\|e_1\|=1$,
\[
 \|b\|^2=\|A(W)-u\|^2+(A_1(W)-u_1)A_1(\mu_S)+\tfrac14A_1(\mu_S)^2 .
\]
Also $\frac14A_1(\mu_S)^2\le\|A(\mu_S)\|^2$, so $\mathcal Q\ge\|a\|^2+\|b\|^2\ge2\|\frac12(a+b)\|^2$. The average of $a$ and $b$ is
$A(\mu_F+\frac12\mu_M)-u+\frac14A_1(\mu_S)e_1$. This gives
\begin{equation}\label{eq:squares}
 \mathcal Q\ge\Delta,\qquad\Delta=2\sum_k\frac1k\bigl(A_k(\mu_F+\tfrac12\mu_M)+\tfrac14A_1(\mu_S)\delta_{k1}-u_k\bigr)^2 .
\end{equation}
The weight $\frac12$ of $\mu_M$ records that these points occur in one of the two Vandermonde products of the $x$ variables, and the weight
$\frac14$ is half of the $\frac12$ in $b$. The diagonal terms are treated by OpenAI's regularization, which replaces the moments by damped moments
with an error that tends to zero, as in the proof of Lemma~\ref{lem:freetangent}. We suppress this.

\subsection{The trial sequences}\label{ssec:trials}
The trial sequences have the form
\begin{equation}\label{eq:trial}
 u_k=c_k+\operatorname{Re}\sum_jr_jz_j^k\qquad(k\ge1;\ c_k=0\text{ for }k>\Kcoef),
\end{equation}
with $\Kcoef$ Chebyshev coefficients and finitely many geometric terms. The geometric terms have a simple
meaning. Since $\sum_{k\ge1}z^kT_k(x)/k=-\frac12\log(1-2xz+z^2)$, a term $(z,r)$ contributes
$-\frac12\operatorname{Re}\,r\log(1-2xz+z^2)$ to the potential $\mathsf T(u,x)$. If
$x_0+i\epsilon_0=\frac12(z+z^{-1})$, then $1-2xz+z^2=2z(x_0+i\epsilon_0-x)$. So the term is, up to a constant,
the logarithmic potential of a point charge of complex strength $r$ at $x_0+i\epsilon_0$ and its conjugate. The
extremal measure has a gap in its support around $x=0$. Its potential has a logarithmic singularity at $x=0$
and square-root behaviour at the edges of the gap. A Chebyshev expansion converges slowly near such points,
while a charge at a small distance $\epsilon_0$ reproduces the local behaviour with one term. This is also why
the quadrature nodes of Section~\ref{sec:real} avoid a neighbourhood of $x=0$. Appendix~\ref{app:field} gives the charges
and the certified bounds.

\subsection{The two cases at the improved rate}\label{ssec:cases}
Theorem~\ref{thm:real} asks for the rate $U_0=U+\Gamma_0/n_0^2$ instead of $U$. In the constrained case the
inequality above holds at $U_0$ for every $\alpha$, with the slope $c_2$. In the case with doubled interaction
Lemma~\ref{lem:energy-ineq} gives it at $U_0$ only for $\alpha\ge\alpha_1$, because the extra term $U_0-U$ is
absorbed by the term $c_2\alpha$ and the quadratic term once $\alpha$ is not small. So a new argument is needed
only in the case with doubled interaction with fewer than $\alpha_1n$ prescribed columns.
Section~\ref{sec:signed} supplies it, with the same constant $c_2$. Precisely, it proves the following.

\begin{proposition}\label{prop:signed-goal}
In the case with doubled interaction, for every configuration with $m\le\alpha_1n$ prescribed columns, the
integrand of Lemma~\ref{lem:andreief}, with the generalized Vandermonde determinant replaced by the bound of
Lemma~\ref{lem:binomvand}, is at most
\[
 \exp\Bigl(n^2\bigl(-U_0+\tfrac12\log2+c_2\alpha+o(1)\bigr)\Bigr)\,C_0^m\prod_j\mathcal B(k_j),\qquad \alpha=m/n,
\]
with an absolute constant $C_0$.
\end{proposition}

\begin{proof}[Proof of Theorem~\ref{thm:real}, given Proposition~\ref{prop:signed-goal}]
By Proposition~\ref{prop:signed-goal} and the discussion above, in both cases a determinant $M_f$ with $m$
prescribed columns $k_1,\ldots,k_m$ is at most
$C_0^m\prod_j\mathcal B(k_j)e^{n_0c_2Nm}\exp\bigl(n_0^2N^2(-U_0+\frac12\log2+o(1))\bigr)$. The factor $C_0^m$ comes
from the expansion of the row functions at the prescribed points. By \eqref{eq:wk} and the bound for
$\mathcal B(k)$, the constants of Lemma~\ref{lem:expansion} satisfy
\[
 \log(B_kC_k)\le\bigl(n_0c_2+w(k/n)+n_0\tilde h_2(k/n)+o(1)\bigr)N\le(\rho+o(1))N .
\]
For $|\varepsilon|\le e^{-\rho N}$ Lemma~\ref{lem:expansion} gives a factor $(1+O(N^5))^n=e^{o(N^2)}$. With
$D=n_0^2(U_0-\frac12\log2)$ this is the claim.
\end{proof}
